\section[Koszul $R$-corings]{Koszul $\boldsymbol{R}$-corings}\label{section4}

In this section we introduce and study the properties of Koszul corings, the dual notion of Koszul $R$-rings. Here we shall also show that any Koszul coring is quadratic.

Strongly graded comodules, the substitute for strongly graded rings when one works with graded corings instead of graded rings, will play an important role in this part of the paper. Hence, we start by discussing their main properties.

Let $(X,\rho^X)$ be a graded right comodule over a graded $R$-coring $C$. For every $n$, the map~$\rho^X$ def\/ines a morphism of graded $C$-comodules $\rho_n^X$ from $X$ to $X_n\ot C$, whose component of degree~$k$ is~$\rho_{n,k-n}^X$. Note that the component of degree $k$ of $X_n\ot C$ is $X_n\ot C_{k-n}$, where $C_i=0$, for all $i<0$. Thus, in particular, $\rho_{n,k-n}^X=0$ for $k<n$.

\begin{Definition}\label{de:strongly_graded_n}
Let $X=\oplus_{n\geq 0} X_n$ be a graded right $C$-comodule, with structure map $\rho^X$. We say that $X$ is \emph{strongly graded in degree} $n$ if the morphism of graded $C$-comodules $\rho^X_{n} \colon X \rightarrow X_n \otimes C$ is injective.
\end{Definition}

Let us remark that, by def\/inition, a graded right $C$-comodule is strongly graded in degree $n$ if and only if $X_p=0$ for all $p<n$ and $\rho_{n,q}\colon X_{n+q}\to X_n\ot C_q$ is injective for all $q\geq 0$.

Some useful properties of graded comodules are collected in the following lemma.

\begin{Lemma}\label{le:strongly_graded_n}
Let $C$ be a strongly graded $R$-coring. Let $(X,\rho^X)$ denote a graded right $C$-comodule.
\begin{enumerate}\itemsep=0pt
 \item[$1.$] The inclusion $\Ker\rho^X_{i,m-i}\subseteq \Ker\rho^X_{j,m-j}$ holds for all $0\leq j\leq i < m$.

 \item[$2.$] There is a canonical isomorphism $\Hom^C(R,X)_m\cong\bigcap_{i=0}^{m-1}\Ker\rho^X_{i,m-i}=\Ker\rho^X_{m-1,1}$.

 \item[$3.$] $X$ is strongly graded in degree $n$ if and only if $X_i=0$, for $i<n$, and $\rho_{i,1}^X$ is injective for $i\geq n$.

 \item[$4.$] Let $f\colon X\to Y$ be a morphism of graded right $C$-comodules. If $X$ is strongly graded in degree $n$ and $f_n\colon X_n\to Y_n$ is an injective map then $f$ is injective as well.
 \end{enumerate}
\end{Lemma}

 \begin{proof}
 To prove (1), consider the following commutative diagram:
 \begin{gather*}
 \xymatrixcolsep{5pc}\xymatrix{ X_m \ar[r]^-{\rho_{i,m-i}^X} \ar[d]_-{\rho_{j,m-j}^X} & X_i \ot C_{m-i} \ar[d]^-{\rho_{j,i-j}^X} \\
X_j \ot C_{m-j} \ar[r]_-{\operatorname{I}_{X_j} \otimes \Delta^C_{i-j,m-i}} & X_j \ot C_{i-j} \ot C_{m-i}.}
 \end{gather*}
Let $x \in \operatorname{Ker}\rho_{i,m-i}^X$. Using the fact that $\operatorname{I}_{X_j} \ot \Delta^C_{i-j,m-i}$ is injective, as $C$ is strongly graded, it follows that $\rho^X_{j,m-j}(x)$ is zero, hence $x \in \operatorname{Ker}\rho_{j,m-j}^X$, as required.

(2) Let $X$ and $Y$ denote two graded $C$-comodule. By def\/inition, $f\in \Hom^C(X,Y)_m$ if and only if $f$ is a morphism of $C$-comodules and $f(X_k)\subseteq Y_{m+k}$, for all $k$. In particular, for a~$C$-colinear map $f\colon R\to X$ of degree $m$, we have $f(1)\in X_m$ and $\rho^X(f(1))=f(1)\ot 1$. The latter relation is equivalent to the equations $\rho^X_{i,m-i}(f(1))=0$, for all $0\leq i\leq m-1$.
In conclusion, the required isomorphism is given by the map $f\mapsto f(1)$. The equation from the statement follows by the f\/irst part of the lemma.

(3) In view of the remark just after the Def\/inition \ref{de:strongly_graded_n}, the condition $X_i=0$ for $i<n$ is part of the hypothesis for both implications. We must prove that $\rho^X_{i,1}$ is injective for all $i\geq n$ if and only if $\rho^X_{n,q} \colon X_{n+q} \to X_n \ot C_q$ is injective for all $q\in\N$.

To prove that $\rho_{n,q}^X$ is injective for $q\in\N$ we proceed by induction on $q$. The map $\rho_{n,0}^M$ coincides with the canonical isomorphism $M_n\cong M_n\ot R$ and $\rho_{n,1}^M$ is injective by assumption. Let us assume that $\rho^X_{n,q}$ is injective. Thus, the horizontal arrow on the bottom of the following commutative diagram is injective,
\begin{gather*}
 \xymatrixcolsep{5pc} \xymatrix{ X_{n+q+1} \ar[r]^-{\rho_{n,q+1}^X} \ar[d]_{\rho_{n+q,1}^X} & X_n \ot C_{q+1} \ar[d]^-{\operatorname{I}_{X_n} \ot \Delta^C_{q,1}} \\
X_{n+q} \ot C_1 \ar[r]_-{\rho^X_{n,q} \ot \operatorname{I}_{C_1}} & X_n \ot C_q \ot C_1.}
\end{gather*}
Since, by hypothesis, the leftmost vertical map is injective it follows that $\rho^X_{n,q+1}$ is also injective.

For the converse one uses the commutative diagram
\begin{gather*}
 \xymatrixcolsep{5pc}\xymatrix{ X_{i+1} \ar[r]^-{\rho^X_{i,1}} \ar[d]_-{\rho^X_{n,i+1-n}} & X_{i} \ot C_1 \ar[d]^-{\rho^X_{n,i-n} \ot \operatorname{I}_{C_1}} \\
X_ n\ot C_{i+1-n} \ar[r]_-{\operatorname{I}_{X_n} \ot \Delta^C_{i-n,1}} & X_n \ot C_{i-n} \ot C_1.}
\end{gather*}
The component $\rho^X_{n,i+1-n}$ is injective by assumption, while $\operatorname{I}_{X_n}\ot\Delta_{i-n,1}^C$ is so as $C$ is strongly graded. We conclude that $\rho^X_{i,1}$ is injective, as required.

(4) We have to prove that all components $f_k$ are injective. For $k<n$ this property trivially holds as, by preceding part of the lemma, we have $X_k=0$. The map $f_n$ is injective by assumption, so we can suppose that $k>n$. The following diagram is commutative, as $f$ is a morphism of graded comodules,
\begin{gather*}
 \xymatrixcolsep{5pc}\xymatrix{ X_{k} \ar[r]^-{f_{k}} \ar[d]_-{\rho^X_{n,k-n}} & Y_{k} \ar[d]^-{\rho^Y_{n,k-n}} \\
X_n \ot C_{k-n} \ar[r]_-{f_n \ot \operatorname{I}_{C_{k-n}}} & Y_n \ot C_{k-n}.}
\end{gather*}
By hypothesis, $f_n\ot\Id_{C_{k-n}}$ is injective. On the other hand, $\rho_{n,k-n}^X$ is injective, as $X$ is strongly graded in degree $n$. We conclude that $x=0$.
 \end{proof}

We can now introduce Koszul corings by dualizing \cite[Def\/inition 1.2.1]{bgs}. Several characterizations of this class of graded corings are given in Theorem \ref{thm:kcoring}.

\begin{Definition} \label{def:kcor}
A connected $R$-coring $C$ is called \emph{Koszul} if $R$ has an resolution $0 \rightarrow R \rightarrow Q^\bullet$ by injective graded right $C$-comodules such that every term $Q^n$ is strongly graded in degree~$n$.
\end{Definition}

\begin{Theorem} \label{thm:kcoring}
Let $C$ be a connected, strongly graded $R$-coring. The following are equivalent:
\begin{enumerate}\itemsep=0pt
\item[$(1)$] the coring $C$ is Koszul;
\item[$(2)$] the pair $(E(C),C)$ is Koszul;
\item[$(3)$] the pair $(C^!,C)$ is Koszul;
\item[$(4)$] the canonical morphism $E(C)\rar C^!$ is an isomorphism;
\item[$(5)$] the $R$-ring $E(C)$ is strongly graded;
\item[$(6)$] the canonical map $QE(C)\to E^1(C)$ is an isomorphism;
\item[$(7)$] the relation $E^{n,m}(C)=0$ holds for all $n \neq m$.
\end{enumerate}
\end{Theorem}

\begin{proof}
Note that, for any Koszul pair $(A,C)$, the complex $\K^\bullet_r(A,C)$ is a resolution of $R$ sa\-tisfying the conditions from the def\/inition of Koszul corings. In particular it follows that $C$ is Koszul, so the implications $(2) \Longrightarrow (1)$ and $(3) \Longrightarrow (1)$ are obvious.

To prove $(1) \Longrightarrow (7)$, assume that $C$ is Koszul. That is, following the def\/inition, $R$ has an injective resolution $0 \rightarrow R \rightarrow Q^\bullet$ such that $Q^n$ is strongly graded in degree $n$. Note that $E^{n,m}(C)=\operatorname{Ext}_{n,m}^C(R,R)$ is the cohomology in degree $n$ of the complex obtained by applying the functor $\operatorname{Hom}^C(R,-)_m$ to this resolution. Thus, to conclude the proof of this implication, it is enough to show that $\operatorname{Hom}^C(R,Q^n)_m$ is trivial for $m\neq n$. By Lemma \ref{le:strongly_graded_n}(2) we have $\operatorname{Hom}^C(R,Q^n)_m=\operatorname{Ker}\rho^{Q^n}_{m-1,1}$. Thus, the claim follows from the fact that $Q^n$ is strongly graded since, in view of Lemma \ref{le:strongly_graded_n}, we have $Q^{n,m}=0$ for $m<n$, and the component $\rho^{Q^n}_{m-1,1}$ is injective for all $m>n$.

For $(7) \Longrightarrow (3)$, let us assume that $E(C)$ is diagonal. We shall prove by induction on $n$ that $\K_r^\bullet(C^!,C) $ is exact in degree $n$. The sequence
\begin{gather*}
 0\longrightarrow R\longrightarrow C_0^!\ot C\longrightarrow
 C_1^!\ot C
\end{gather*}
is exact by the def\/inition of $\K^\bullet(C^!,C)$ and the fact that $C$ is strongly graded, cf.\ \cite[Lemma~2.1]{jps}.


We now assume that $\K^\bullet(C^!,C)$ is exact in degree $i$, for all $0\leq i\leq n-1$, and prove that it is exact in degree $n$ as well. We consider the exact sequence
\begin{gather}\label{eq:sequenceC!}
 0 \rightarrow R \rightarrow C^!_0\ot C \xrightarrow{d^0_r} C_1^! \otimes C \xrightarrow{d^1_r} \cdots \xrightarrow{d^{n-2}_r} C^!_{n-1}\ot C \xrightarrow{d^{n-1}_r} C_n^! \otimes C \xrightarrow{\al}Y \rightarrow 0,
\end{gather}
where $Y:=\Coker d^{n-1}_r$ and $\al$ denotes the canonical projection. We claim that $Y$ is strongly graded in degree $n+1$. Since $d_r^{n-1}(\widehat{x}\ot c)= \widehat{x\ot c}\ot 1$, for all $x\in C_1^{n-1}$ and $c\in C_1$, the homogeneous component $(C_n^!\ot C)_n$ is included into the image of $d_r^{n-1}$. Thus $Y$ has no non-zero elements of degree $n$. On the other hand, if $Y_k$ denotes $k$-degree component of $Y$, then $Y_k=0$ for all $k<n$, as $Y$ is the quotient of $C_n^!\ot C$ by a graded subcomodule. In conclusion, for proving our claim, it remains to show that $\rho^Y_{m-1,1}$ is injective for all $m>n+1$.

Let $X:=\Ker \al=\im d^{n-1}_r$. We complete $0 \rightarrow {X} \rightarrow C_n^! \otimes C$ to a resolution
\begin{gather*}
 0 \rightarrow X\rightarrow C_n^!\ot C \rightarrow Q^{n+1} \rightarrow Q^{n+2} \rightarrow \cdots
\end{gather*}
by injective graded right $C$-comodules, which combined with \eqref{eq:sequenceC!} yields us a new injective resolution
\begin{gather*} %\label{resinj}
0 \rightarrow R \rightarrow C_0^!\ot C \rightarrow C_1^!\ot C\rightarrow \dots \rightarrow C_n^! \otimes C \rightarrow Q^{n+1} \rightarrow Q^{n+2} \rightarrow \cdots.
\end{gather*}
Thus we can compute $\operatorname{Ext}_C^{n+1,m}(R,R)$ as the cohomology in degree $n+1$ of the complex:
\begin{gather*}
 0 \rightarrow \operatorname{Hom}^C(R,C_0^!\ot C)_m \rightarrow \operatorname{Hom}^C(R,C_1^!\ot C)_m \rightarrow \dots \rightarrow \operatorname{Hom}^C(R, C_n^!\ot C)_m\\
 \hphantom{0}{} \rightarrow \operatorname{Hom}^C\big(R,Q^{n+1}\big)_m \rightarrow \cdots.
\end{gather*}
Since $\operatorname{Ext}^{1,m}_C(R,X)$ is the $1$-st cohomology group of the complex
\begin{gather*}
 0 \rightarrow \operatorname{Hom}^C(R,C_n^! \otimes C)_m \rightarrow \operatorname{Hom}^C(R,Q^{n+1})_m \rightarrow \operatorname{Hom}^C\big(R,Q^{n+2}\big)_m \rightarrow \cdots
\end{gather*}
we deduce that $\operatorname{Ext}^{1,m}_C(R,X)\cong\operatorname{Ext}_C^{n+1,m}(R,R)$. Recall that $X$ is the kernel of $\al$, so the long exact sequence connecting the functors of $\Ext^{\bullet,m}_C(R,-)_m$ can be written as follows
\begin{gather*}
 0 \rightarrow \operatorname{Hom}^C(R,X)_m \rightarrow \operatorname{Hom}^C(R, C_n^! \otimes C)_m \rightarrow \operatorname{Hom}^C(R,Y)_m \rightarrow \operatorname{Ext}^{1,m}_C(R,X) \rightarrow \cdots.
\end{gather*}
From the standing assumption, $E(C)$ is diagonal, so using the above isomorphism between the Ext-groups, it follows that $\operatorname{Ext}_C^{1,m}(R,X)=0$, for all $m>n+1$. Moreover, $\operatorname{Hom}^C(R,C_n^! \otimes C)$ is isomorphic to $\operatorname{Hom}_R(R,C_n^! )$ as graded $R$-modules, where the latter module is concentrated in degree~$n$. It follows that $\operatorname{Hom}^C(R,Y)_m=0$, for all $m > n+1$. By Lemma~\ref{le:strongly_graded_n}(2) it follows that $\rho^Y_{m-1,1}$ is injective for all $m>n+1$. Since we already know that $Y_k=0$ for $k<n+1$, we conclude by Lemma~\ref{le:strongly_graded_n}(3) that $Y$ is strongly graded in degree $n+1$, that is our claim has been proved.

We can now prove that $\K^\bullet(C^!,C)$ is exact in degree $n$. We must show that the kernel of the map $\pa\colon Y\to C^!_{n+1}\ot C$ induced by $d^n_r$ is trivial. In view of Lemma~\ref{le:strongly_graded_n}(4), as $Y$ is strongly graded in $n+1$, it is enough to prove that the component $\pa_{n+1}\colon Y_{n+1}\to C_{n+1}\ot C_0$ of $\pa$ is injective.

Let $y$ be an element in the kernel of $\pa_{n+1}$. Hence there are ${y}_1, \dots,{y}_p\in C^{(n)}$ and $c_1,\dots, c_p\in C_1$ such that $y$ is the equivalence class of $\sum\limits_{i=1}^p \widehat{y}_i\ot c_i$ in $Y_{n+1}=(C_n^!\ot C_1)/d_r^{n-1}(C_{n-1}^!\ot C_2)$. Here $\widehat{y}_i$ denotes the class of $y_i$ in $C^!_n=C_1^{(n)}/W_n$, where $W_n=\sum\limits_{k=1}^{n-1} C_1^{(k-1)}\ot\im\Delta_{1,1}\ot C_1^{(n-k-1)}$. Since $\pa_{n+1}(y)=0$ and using the fact that $d_r^n$ is induced by the multiplication in $C^!$, it follows that $y'=\sum\limits_{i=1}^p {y}_i\ot c_i$ belongs to the submodule $W_{n+1}\subseteq C^{(n+1)}$ from the construction of $C_{n+1}^!$. Note that $W_2=\im\Delta_{1,1}$. Thus $y'=x+\sum\limits_{j=1}^q y_j''\ot c^j_{1,1}\ot c^j_{2,1}$, for some $x\in W_n\ot C_1$, $y''_1,\dots,y''_q\in C_1^{(n-1)}$ and $c^1,\dots,c^q\in C_2$. The relation $y=0$ now follows by the computation
\begin{gather*}
 \sum_{i=1}^p \widehat{y}_i\ot c_i=\sum_{j=1}^q \widehat{y_j''\ot c^j_{1,1}}\ot c^j_{2,1}=d_r^{n}\left(\sum_{j=1}^q {\widehat{y_j''}\ot c^j}\right).
\end{gather*}

For the implication $(3) \Longrightarrow (4)$ we use the morphism of complexes $\phi^\bullet\colon \beta^\bullet_r(C) \rar \operatorname{K}_r^\bullet(C^!,C)$, which lifts the identity of $R$, see Section~\ref{fa:exemples_morphisms}. Since both $\K^{\bullet}_{r}(C^!,C)$ and $\beta{}^{\bullet}_{r}(C)$ are resolutions of $R$ by injective right $C$-comodules, by the comparison theorem, the morphism $\{\phi^n\}_{n\in\N}$ is invertible up to a homotopy in the category of complexes of right $C$-comodules. Hence $\Hom^C(R,\phi^\bullet)$ is invertible up to a homotopy in the category of complexes of right $R$-modules. It follows that $\h^n(\Hom^C(R,\phi^\bullet))$ is an isomorphism for all $n\geq 0$. Thus $\phi^n_C\colon E_n(C)\to C^!_n$ is an isomorphism, as it coincides with $\h^n(\Hom^C(R,\phi^\bullet))$, see Section~\ref{fa:exemples_morphisms}.

The implication $(4) \Longrightarrow (5)$ follows immediately since $C^!$ is always strongly graded, hence $E(C)$ is so. Furthermore, $(5) \iff (6)$ and $(5) \Longrightarrow (7)$ are direct consequences of Lemma~\ref{lema:ring}.

We conclude the proof by remarking, for the implication $(3) \Longrightarrow (2)$, that the complex $\K^\bullet_r(C^!,C)$ is isomorphic to $\K^\bullet_r(E(C),C)$, which makes the latter exact.
\end{proof}
